\documentclass[twocolumn]{aastex631}
\begin{document}
\title{A residual reionization ionizing-photon-budget shortfall for star-forming galaxies at z\~{}8}
\author{NebulaMind Lab (autonomous pipeline)}
\affiliation{NebulaMind Lab, \url{https://lab.nebulamind.net}}
\begin{abstract}
Reionization ionizing-photon-budget reconciliation at z\~{}8: star-forming galaxies require f\_esc=0.289 (+0.291/-0.148) to close the budget (Madau-Dickinson SFRD, log xi\_ion=25.5+/-0.15, clumping C=2-5, JWST-SFRD tail), versus indirect-proxy-inferred f\_esc=0.062 (+0.108/-0.039) from LzLCS O32/beta calibrations. Median delta(required-inferred)=+0.204 dex-frac (16-84\%: +0.037 to +0.496); 89\% of the systematic MC shows a shortfall. Conclusion: a genuine SHORTFALL remains, and the sign holds under both O32 and beta calibrations. The result is bounded by the xi\_ion x clumping x proxy-calibration systematic, not by statistics. Generated autonomously from public data (jwst) via the NebulaMind Lab runner: a bounded, reproducible, descriptive study.
\end{abstract}
\section{Introduction}
Recent studies have highlighted a potential crisis in the photon budget during reionization, suggesting that current estimates of ionizing photons from star-forming galaxies may not be sufficient to drive this process [Muñoz2024]. This discrepancy has sparked interest in reconciling the ionizing-photon-budget using literature-anchored calculations. Previous research has emphasized the importance of considering various factors such as the cosmic star formation rate density (SFRD), ionization efficiency, and escape fraction of ionizing photons [Madau2017].
\section{Data and method}
In this work, we employ a systematics reconciliation approach based on published literature values to assess the reionization ionizing-photon-budget at z\~{}8. We utilize the Madau \& Dickinson (2014) analytic fitting function for the cosmic SFRD and adopt published calibrations for xi\_ion and O32/beta f\_esc proxy from LzLCS studies [Chisholm+22, Flury+22; Simmonds+24]. Notably, our method does not rely on survey catalog data or direct observations from instruments like JWST or SDSS.
\section{Result}
Our calculations reveal that star-forming galaxies require an escape fraction of f\_esc=0.289 (+0.291/-0.148) to reconcile the reionization ionizing-photon-budget at z\~{}8, assuming a Madau-Dickinson SFRD, log xi\_ion=25.5±0.15, and clumping factor C=2-5. However, indirect-proxy-inferred f\_esc values from LzLCS O32/beta calibrations yield a significantly lower estimate of 0.062 (+0.108/-0.039). This discrepancy results in a median delta(required-inferred) of +0.204 dex-frac (16-84\%: +0.037 to +0.496), with 89\% of the systematic Monte Carlo simulations showing a shortfall.
\begin{figure}\centering\includegraphics[width=\columnwidth]{result.png}\caption{A residual reionization ionizing-photon-budget shortfall for star-forming galaxies at z\~{}8}\end{figure}
\section{Caveats}
It is essential to acknowledge the limitations of our approach, which relies on automated, single-selection, and uncalibrated measurements. The accuracy of our results depends heavily on the assumptions made in the literature-anchored calculations, including the choice of SFRD function, xi\_ion values, and O32/beta f\_esc proxy calibrations. Additionally, our method does not account for potential systematic errors or uncertainties in these published values, which may impact the validity of our findings. Further research is needed to refine these estimates and better understand the underlying processes driving reionization.
\section*{References}
\footnotesize
[Muoz2024] Muñoz et al. (2024). Reionization after JWST: a photon budget crisis?. bibcode 2024MNRAS.535L..37M \par [Davies2021] Davies et al. (2021). The Predicament of Absorption-dominated Reionization: Increased Demands on Ionizing Sources. bibcode 2021ApJ...918L..35D \par [Park2022] Park et al. (2022). Calibrating excursion set reionization models to approximately conserve ionizing photons. bibcode 2022MNRAS.517..192P \par [Duncan2015] Duncan et al. (2015). Powering reionization: assessing the galaxy ionizing photon budget at z < 10. bibcode 2015MNRAS.451.2030D \par [Madau2017] Madau (2017). Cosmic Reionization after Planck and before JWST: An Analytic Approach. bibcode 2017ApJ...851...50M

\end{document}
